\chapter{Постановка задачи и исходная архитектура Otterbrix}

\section{Общая характеристика проекта}

Otterbrix~--- встраиваемый вычислительный фреймворк с открытым
исходным кодом, предназначенный для обработки полуструктурированных
данных в сочетании с классической реляционной моделью. Проект
развивается как альтернатива специализированным документным СУБД
(MongoDB, Couchbase) и аналитическим колоночным системам
(ClickHouse, DuckDB), предлагая унифицированную модель, в которой
SQL-запросы могут одновременно обращаться к структурированным
таблицам и к полям JSON-документов в рамках одного исполнителя.

Архитектурно Otterbrix построен на следующих принципах современных
СУБД:
\begin{itemize}
    \item \textbf{Акторная модель параллелизма} (реализована поверх
    библиотеки actor-zeta)~--- все ключевые компоненты (дисковый
    менеджер, WAL, диспетчер запросов, менеджер индексов)
    оформлены как независимые актёры, обменивающиеся сообщениями.
    Это даёт предсказуемую изоляцию состояния и позволяет
    масштабировать систему на большое число ядер без глобальных
    блокировок.
    \item \textbf{Корутины C++20} (co\_await~/~co\_return)~---
    используются для организации асинхронного исполнения планов
    запросов. Каждый план может приостанавливаться на ожидании
    диска или других актёров, не блокируя нить пула.
    \item \textbf{Полиморфные аллокаторы} (\texttt{std::pmr})~---
    вся работа с памятью идёт через конкретный
    \texttt{memory\_resource}, что позволяет локализовать
    аллокации запросов и контролировать фрагментацию.
    \item \textbf{Собственный SQL-парсер на базе
    libpg\_query}~--- грамматика PostgreSQL обеспечивает
    совместимость с широким спектром существующих запросов и
    инструментов.
    \item \textbf{Колоночное хранилище}~--- таблицы организованы
    как последовательность row~groups, каждый из которых содержит
    независимые колоночные сегменты. Это стандартный приём для
    ускорения аналитических запросов.
\end{itemize}

На момент начала работы проект находился в активной фазе
рефакторинга: команда завершала переход с гибридной модели
хранения (отдельные подсистемы для документов и таблиц) на единое
колоночное хранилище. Этот переход, оформленный одним крупным
коммитом (\texttt{a34d3823}, <<Remove components/document and
migrate to columnar data\_chunk storage>>, февраль 2026), удалил
более 18~000 строк кода и послужил отправной точкой для настоящей
работы.

\section{Двойственность представления: документы и таблицы}

До указанного рефакторинга Otterbrix поддерживал два
принципиально различных типа коллекций.

\textbf{Документные коллекции} использовали отдельный компонент
\texttt{components/document/}, реализовавший собственный вариант
иерархического JSON-подобного значения. Каждый документ
представлял собой дерево (объекты, массивы, скалярные листья),
сериализованное в формате, близком к BSON. Доступ к полям
происходил по JSON-указателю (\texttt{/commit/record/text}),
каждый документ хранился как неделимое значение.

\textbf{Табличные коллекции} использовали колоночное хранилище
\texttt{components/table/} с фиксированной схемой (набор колонок
фиксированных типов). Доступ к ячейкам происходил по паре
(row\_id, column\_id); каждая колонка хранилась в отдельном
сегменте памяти или диска.

Обе реализации имели собственную подсистему хранения, собственные
операторы исполнения (match, group, sort, insert, update, delete),
собственные адаптеры диска, собственный WAL-формат. Класс
\texttt{document\_t} (\texttt{components/document/document.hpp},
удалён в \texttt{a34d3823}) насчитывал около 1100~строк кода и
реализовывал полнофункциональный JSON-документ с поддержкой
мутаций через \texttt{set(json\_pointer, value)},
\texttt{remove(json\_pointer)}, \texttt{set\_array},
\texttt{set\_dict}, а также операций сравнения и сериализации.

Основные недостатки документной реализации:
\begin{itemize}
    \item \textbf{Неэффективный ввод-вывод при аналитических
    запросах.} Для ответа на запрос вида \texttt{SELECT COUNT(*)
    FROM t WHERE commit.operation = 'create'} система вынуждена
    была читать с диска полный документ целиком, хотя для ответа
    требовалось единственное поле.
    \item \textbf{Низкая локальность кэша L1/L2.} Каждое
    обращение по JSON-указателю представляло собой обход дерева
    узлов, расположенных в куче и связанных указателями.
    \item \textbf{Невозможность применения векторных
    оптимизаций.} Векторизованные операторы, повсеместно
    используемые в современных колоночных СУБД (пакетное
    сравнение, SIMD), требуют компактного, однотипного
    расположения значений в памяти.
    \item \textbf{Удвоенная кодовая база.} Наличие двух
    параллельных реализаций операторов приводило к постоянной
    рассинхронизации: оптимизации одного пути не переносились в
    другой; баги всплывали в двух местах независимо.
\end{itemize}

\section{Миграция на единое колоночное хранилище}

Коммит \texttt{a34d3823} завершил миграцию: компонент
\texttt{components/document/} был удалён полностью, все коллекции
стали хранить данные в формате \texttt{data\_chunk\_t}~---
структуры из \texttt{components/vector/}, представляющей собой
вектор колонок с общим числом строк. Каждая колонка хранится в
виде сегментов \texttt{vector\_t}, поддерживающих несколько
представлений (flat, constant, dictionary) и плоский массив
значений одного типа.

После миграции таблицы получили два подвида:
\begin{itemize}
    \item \textbf{Таблицы с фиксированной схемой.} Колонки
    задаются при \texttt{CREATE TABLE db.t (id INT, name TEXT);}
    и остаются неизменными в течение жизни таблицы.
    \item \textbf{Таблицы с вычисляемой (computed) схемой.}
    Создаются без предварительного перечисления колонок~---
    \texttt{CREATE TABLE db.t ();}. Каждая операция
    \textsc{insert} способна добавить новые поля; схема
    расширяется по мере поступления данных. Именно этот режим
    заменил бывшие <<документные>> коллекции и стал основным
    предметом настоящей работы.
\end{itemize}

Миграция закрыла проблему удвоенной кодовой базы, но оставила
открытым ключевой вопрос: \textbf{как именно хранить колонки,
заполненные лишь в небольшой доле строк}. Прямолинейный вариант~---
создавать физическую колонку на каждый путь JSON и заполнять её
NULL для отсутствующих значений~--- порождает две новые проблемы:
\begin{enumerate}
    \item \textbf{Катастрофический рост пространства.} Реальные
    полуструктурированные данные (в том числе стандартный датасет
    JSONBench, сформированный из публичного потока сети Bluesky)
    содержат десятки тысяч различных JSON-путей, большинство из
    которых встречаются в менее чем 0.1\% документов.
    \item \textbf{Замедление записи.} Каждое добавление новой
    колонки в существующую таблицу требует выделения буфера для
    всех $N$ ранее вставленных строк и заполнения его значениями
    по умолчанию.
\end{enumerate}

\section{Формулировка задачи}

Цель выпускной квалификационной работы сформулирована следующим
образом: спроектировать и реализовать эффективный механизм
хранения полуструктурированных данных в рамках колоночного
движка Otterbrix, сочетающий производительность колоночного
доступа к часто встречающимся полям с компактным представлением
редко встречающихся полей, без необходимости поддерживать
отдельную подсистему для JSON-документов.

Для достижения поставленной цели решаются следующие задачи:
\begin{enumerate}
    \item Проанализировать существующие подходы к хранению
    полуструктурированных данных (BSON в PostgreSQL, EAV,
    flatten-представление путей, тип \texttt{JSON} в ClickHouse).
    \item Спроектировать гибридное представление, сочетающее
    колоночное хранение плотных полей и компактное хранение
    редких полей, с автоматической миграцией между режимами.
    \item Реализовать предложенное представление в кодовой базе
    Otterbrix, встраиваясь в существующие механизмы SQL-парсера,
    диспетчера и подсистемы хранения.
    \item Обеспечить корректное поведение системы при
    параллельных операциях \textsc{insert} и \textsc{select} над
    таблицами с вычисляемой схемой.
    \item Выполнить оптимизацию физических операторов (проекция
    колонок, типизированные предикаты, векторизация) для
    достижения производительности, сопоставимой с существующими
    промышленными решениями.
    \item Провести сравнительное тестирование на бенчмарке
    JSONBench с ClickHouse и PostgreSQL в качестве референсных
    систем.
\end{enumerate}

\section{Научная новизна и практическая значимость}

Основная новизна работы заключается в предложении механизма
\textbf{гетерогенного хранения колонок}: в рамках одной логической
таблицы часть колонок хранится в основной таблице (колоночное
хранение), часть~--- в отдельных вспомогательных таблицах формата
\texttt{(\_id, value)} (разреженное хранение), с автоматической
миграцией полей между режимами по мере накопления данных. В отличие
от подхода ClickHouse, где редкие поля попадают во внутренний
<<shared~data>> формат с ограниченной функциональностью,
предложенное решение сохраняет для разреженных колонок полный
набор операций (произвольные предикаты, индексы, агрегации) за счёт
реализации чтения через LEFT~JOIN на уровне исполнителя.

Практическая значимость состоит в получении системы, которая
обходит известные ограничения наивного колоночного хранения JSON
(неограниченный рост схемы, раздувание нулевыми значениями), не
требует поддержки отдельной подсистемы для документов, демонстрирует
производительность, сопоставимую с ClickHouse и существенно
превосходящую PostgreSQL с JSONB, и реализована на основе
существующих в Otterbrix механизмов колоночного хранения, что
позволяет переиспользовать все накопленные оптимизации
(векторизация, zone~maps, projection~pushdown) без дополнительных
усилий.

\chapterconclusion

В настоящей главе описана исходная архитектура проекта Otterbrix,
мотивация отказа от документной подсистемы хранения в пользу
единой колоночной архитектуры, и сформулированы задачи, решаемые
в настоящей выпускной квалификационной работе. Последующие главы
посвящены обзору существующих подходов, проектированию и реализации
предлагаемого решения.
